\subsection{Model assumptions}

The \textbf{main} assumptions used to model the steam turbine cycle are the following :
\begin{itemize}
    \setlength\itemsep{0em}
    \item The combustion process is modelled as a simple adiabatic heat transfer.
    \item The pressure losses in the steam generator are taken into account.
    \item Isobaric condensation.
    \item Adiabatic expansions in the turbines (the isentropic model is used).
    \item The pumps are electrically powered rather than mechanically driven by the same shaft as the turbines.
    \item For the exergy analysis, the expansion valves between each feedwater heater are neglected.
\end{itemize}

\subsection{Methodology}

In order to develop the model, we first computed the thermodynamic quantities for all states of the cycle.
Based on these quantities, we calculated the different flow rates
in the bleedings through energy balances at each feedwater heater.
The total flow rate of the cycle was determined using the power rating.
Finally, the energy and exergy losses were calculated using energy and 
exergy balances on each component of the cycle,
and the energetic and exergetic efficiencies were determined.
